\section{Numerical Experiments}
\label{sec:experiments}

Non-uniform global scheduling should improve sample quality whenever per-timestep
noise sensitivity varies, a condition our theory predicts and offline profiling
confirms.
We test this prediction across multiple diffusion model families under strict
NFE budgets, varying architectures, prompt sets, and local search operators.

\subsection{Experimental Setting}
\label{sec:exp-setting}
We evaluate GAINS on both Stable Diffusion and EDM under \emph{strict} NFE budgets.
Two lightweight verifiers, adopted as single-image rewards in prior noise-trajectory-search benchmarks, serve as objectives: \textbf{Brightness}, computed from perceived luminance, and \textbf{Compressibility}, derived from the maximally-compressed JPEG file size and normalized to a fixed range.%
\footnote{Because resolution is fixed in our setup, JPEG size is determined by image compressibility.}
Both verifiers enable controlled, repeatable measurements of test-time scaling without human annotation.

The \textbf{benchmark method} (``Naive'') pairs an $\epsilon$-greedy noise search operator with a \textbf{uniform per-timestep} compute schedule, following prior work.
GAINS retains the same local search operators but replaces the uniform policy with a two-stage scheduler: offline profiling yields a coarse per-step allocation, and online control applies windowed early stopping driven by (i) observed score gains and (ii) candidate score variance.
Early stopping acts only within the high-value region identified by offline profiling; saved NFE is redistributed to the first low-value steps so that total NFE exactly matches the prescribed budget.
All experiments run on a single NVIDIA A100 GPU; wall-clock time is dominated by UNet evaluations and therefore determined by the NFE budget.

\subsection{Stable Diffusion: Scaling with NFE}
\label{sec:exp-sd-scaling}
\begin{table}[t]
\centering
\caption{SD results across different NFE budgets.}
\label{tab:sd_budget}
\begin{tabular}{c|c|c}
\toprule
\textbf{NFE} & \textbf{Naive} & \textbf{GAINS} \\
\midrule
\multicolumn{3}{c}{\textit{Brightness}} \\
\midrule
400 & $0.7025 \pm 0.0466$ & $0.7248 \pm 0.0563$ \\
500 & $0.7094 \pm 0.0631$ & $0.7346 \pm 0.0599$ \\
800 & $0.7511 \pm 0.0620$ & $0.7832 \pm 0.0462$ \\
\midrule
\multicolumn{3}{c}{\textit{Compressibility}} \\
\midrule
400 & $0.8833 \pm 0.0166$ & $0.8946 \pm 0.0164$ \\
500 & $0.8825 \pm 0.0170$ & $0.9004 \pm 0.0147$ \\
800 & $0.8892 \pm 0.0198$ & $0.9008 \pm 0.0123$ \\
\bottomrule
\end{tabular}
\end{table}

GAINS improves both absolute scores and the \emph{rate} at which quality scales with NFE on Stable Diffusion (\cref{tab:sd_budget}).
For \textit{Brightness}, margins over uniform allocation widen from +0.0223 (400 NFE)
to +0.0252 (500 NFE) and +0.0321 (800 NFE).
The GAINS 400-NFE result ($0.7248$) already surpasses uniform allocation at
500 NFE ($0.7094$), reaching higher brightness with
\textbf{20\% fewer} function evaluations.
The advantage is more pronounced for \textit{Compressibility}:
uniform-allocation scores stagnate from $0.8833$ (400) to $0.8825$ (500) and $0.8892$ (800),
whereas GAINS reaches $0.8946$, $0.9004$, and $0.9008$.
GAINS at \textbf{400 NFE} ($0.8946$) already exceeds uniform allocation
at \textbf{800 NFE} ($0.8892$), achieving \textbf{50\% fewer} function evaluations for
higher compressibility.
This consistent leftward shift of the quality curve confirms that
per-step-aware scheduling converts NFE into quality more efficiently than uniform allocation.

\subsection{EDM: Scaling with NFE}
\label{sec:exp-edm-scaling}
\begin{table}[t]
\centering
\caption{EDM results across different NFE budgets.}
\label{tab:edm_budget}
\begin{tabular}{c|c|c}
\toprule
\textbf{NFE} & \textbf{Naive} & \textbf{GAINS} \\
\midrule
\multicolumn{3}{c}{\textit{Brightness}} \\
\midrule
144 & $0.9507 \pm 0.0153$ & $0.9887 \pm 0.0046$ \\
180 & $0.9617 \pm 0.0160$ & $0.9904 \pm 0.0036$ \\
288 & $0.9787 \pm 0.0178$ & $0.9988 \pm 0.0012$ \\
\midrule
\multicolumn{3}{c}{\textit{Compressibility}} \\
\midrule
144 & $0.6714 \pm 0.0165$ & $0.6845 \pm 0.0074$ \\
180 & $0.6774 \pm 0.0120$ & $0.7003 \pm 0.0089$ \\
288 & $0.6901 \pm 0.0118$ & $0.7165 \pm 0.0085$ \\
\bottomrule
\end{tabular}
\end{table}
EDM exhibits an even stronger acceleration (\cref{tab:edm_budget}).
For \textit{Brightness}, GAINS achieves
$0.9887$, $0.9904$, and $0.9988$ at 144, 180, and 288 NFE, corresponding to gaps of +0.0380, +0.0287,
and +0.0201 over uniform allocation.
The GAINS 144-NFE brightness ($0.9887$) already surpasses uniform allocation at
288 NFE ($0.9787$), delivering higher quality with \textbf{50\% fewer}
function evaluations.
Similarly, at 180 NFE GAINS ($0.9904$) exceeds uniform allocation at 288 NFE,
a \textbf{37.5\%} compute reduction.
This pattern is consistent with \cref{prop:offline}(iii): EDM's concentrated
mid-trajectory sensitivity creates large variance dispersion $\{\sigma_t\}$,
amplifying the benefit of non-uniform allocation.
For \textit{Compressibility}, GAINS reaches $0.6845$, $0.7003$, and $0.7165$
versus uniform-allocation scores of $0.6714$, $0.6774$, and $0.6901$.
The GAINS 144-NFE result ($0.6845$) already exceeds uniform allocation at 180 NFE
($0.6774$), saving \textbf{20\%} NFE;
the GAINS 180-NFE result ($0.7003$) surpasses uniform allocation at 288 NFE ($0.6901$),
saving \textbf{37.5\%} NFE.
Notably, improvements grow from +0.0131 to +0.0229 and
+0.0264 as budget increases, indicating that global scheduling allocates
incremental NFE to impactful timesteps rather than spreading it uniformly.

\subsection{Ablation: Offline Allocation vs.\ Offline+Online Control}
\label{sec:exp-ablation-offline-online}
\begin{table}[t]
\centering
\caption{Stable Diffusion results under fixed NFE $=400$.
\textbf{Naive} uses uniform per-step allocation.
\textbf{GAINS} uses offline profiling plus online control.
\textbf{Offline Only} removes online control and uses only offline allocation.}
\label{tab:sd_400_ablation}
\begin{tabular}{lccc}
\toprule
\textbf{Metric} & \textbf{Naive} & \textbf{GAINS} & \textbf{Offline Only} \\
\midrule
Brightness        & $0.7025 \pm 0.0466$ & $\mathbf{0.7248 \pm 0.0563}$ & $0.7158 \pm 0.0645$ \\
Compressibility   & $0.8833 \pm 0.0166$ & $\mathbf{0.8946 \pm 0.0164}$ & $0.8873 \pm 0.0183$ \\
\bottomrule
\end{tabular}
\end{table}

Separating the two scheduling stages reveals complementary contributions (\cref{tab:sd_400_ablation}).
\textbf{Offline Only} already improves over uniform allocation on both metrics, confirming that coarse per-step budgeting is beneficial by itself.
Adding online feedback and windowed early stopping yields the highest scores for both \textit{Brightness} (0.7248) and \textit{Compressibility} (0.8946).
Offline profiling thus captures \emph{where} search matters, while online control refines \emph{how much} budget is worth spending on a given prompt and timestep, matching the two-factor decomposition in \cref{prop:online}(ii).

\subsection{Larger Prompt Set Evaluation}
\label{sec:exp-larger-prompts}
\begin{table}[t]
\centering
\caption{Larger prompt set evaluation.
We use 20 prompts, each repeated 10 times.
Baseline corresponds to the uniform allocation method; GAINS is the proposed global scheduler.}
\label{tab:larger_prompt_exp}
\begin{tabular}{lcc}
\toprule
 & \textbf{brightness} & \textbf{compressibility} \\
\midrule
Baseline          & $0.6949 \pm 0.0756$ & $0.7938 \pm 0.0892$ \\
GAINS             & $\mathbf{0.7213 \pm 0.0752}$ & $\mathbf{0.8076 \pm 0.0962}$ \\
\bottomrule
\end{tabular}
\end{table}

A fixed offline profile might overfit to a narrow calibration set. To test generalization, we expand to 20 prompts with 10 repeats each (\cref{tab:larger_prompt_exp}).
GAINS raises mean brightness from 0.6949 to 0.7213 and mean compressibility from 0.7938 to 0.8076.
Because improvements persist under repeated sampling across a diverse prompt set, the benefit reflects systematic compute reallocation to high-impact timesteps rather than cherry-picked prompt selection.

\subsection{Compatibility with Different Local Search Operators}
\label{sec:exp-local-operator}
\begin{table}[t]
\centering
\caption{Global scheduling combined with different local search operators.
Left block uses uniform allocation; right block uses GAINS.
B: Brightness, C: Compressibility.}
\label{tab:local_search_combo}
\begin{tabular}{lcccc}
\toprule
& \multicolumn{2}{c}{\textbf{Naive}} & \multicolumn{2}{c}{\textbf{GAINS}} \\
\cmidrule(lr){2-3} \cmidrule(lr){4-5}
\textbf{Metric} & \textbf{Zero-order} & \textbf{Random} & \textbf{Zero-order} & \textbf{Random} \\
\midrule
\textbf{B} & $0.4920 \pm 0.0723$ & $0.7382 \pm 0.0392$ & $\mathbf{0.5080 \pm 0.0752}$ & $\mathbf{0.7457 \pm 0.0452}$ \\
\textbf{C} & $0.5727 \pm 0.0660$ & $0.8417 \pm 0.0268$ & $\mathbf{0.5832 \pm 0.0613}$ & $\mathbf{0.8454 \pm 0.0252}$ \\
\bottomrule
\end{tabular}
\end{table}

Because GAINS schedules \emph{when} to search rather than \emph{how}, it should compose with any local search operator.
\cref{tab:local_search_combo} confirms this: under both \textbf{Zero-order} and \textbf{Random} local search, replacing uniform allocation with GAINS improves \textit{Brightness} and \textit{Compressibility}.
The scheduler targets where to spend compute and pairs with different per-timestep proposal mechanisms without modification.

\subsection{Applicability to Flow-based Models}
\label{sec:exp-flow}

Flow-based models pose a distinct challenge: their deterministic ODE transitions
contain no noise to search over.
Following \citet{dockhorn2024stochastic} and \citet{kim2025rbf}, we inject
stochasticity at selected timesteps, converting the ODE into an SDE with the
same marginal distributions.
Once cast in the form of \cref{eq:transition}, GAINS applies without modifying the pretrained model.

\begin{table}[t]
\centering
\caption{Flow-based model results. Stochasticity is introduced via noise
injection~\citep{dockhorn2024stochastic}, after which GAINS
allocates search compute non-uniformly across timesteps.}
\label{tab:flow_results}
\begin{tabular}{c|c|c}
\toprule
\textbf{NFE} & \textbf{Naive} & \textbf{GAINS} \\
\midrule
\multicolumn{3}{c}{\textit{Brightness}} \\
\midrule
80  & $0.5492 \pm 0.0032$ & $\mathbf{0.5507 \pm 0.0028}$ \\
100 & $0.5499 \pm 0.0033$ & $\mathbf{0.5519 \pm 0.0033}$ \\
150 & $0.5515 \pm 0.0039$ & $\mathbf{0.5526 \pm 0.0028}$ \\
\midrule
\multicolumn{3}{c}{\textit{Compressibility}} \\
\midrule
80  & $0.8135 \pm 0.0050$ & $\mathbf{0.8162 \pm 0.0031}$ \\
100 & $0.8156 \pm 0.0044$ & $\mathbf{0.8170 \pm 0.0063}$ \\
150 & $0.8157 \pm 0.0058$ & $\mathbf{0.8193 \pm 0.0046}$ \\
\bottomrule
\end{tabular}
\end{table}

Global scheduling improves quality even under the reduced stochasticity
of flow models (\cref{tab:flow_results}).
For \textit{Brightness}, GAINS yields gains of
+0.0015, +0.0021, and +0.0011 at 80, 100, and 150 NFE;
the 80-NFE result ($0.5507$) already matches
uniform allocation at 100 NFE ($0.5499$), saving \textbf{20\%} NFE.
For \textit{Compressibility}, improvements reach +0.0027, +0.0014,
and +0.0036;
the GAINS 100-NFE score ($0.8170$) exceeds uniform allocation at 150 NFE
($0.8157$), a \textbf{33\%} compute reduction.
Margins are smaller than on SD and EDM, consistent with the narrower
variance dispersion $\{\sigma_t\}$ after ODE-to-SDE conversion
(\cref{prop:offline}(iii) predicts that gains shrink as per-step variances become
more uniform).
Nonetheless, the quality curve shifts leftward at every budget,
and GAINS also reduces run-to-run variance.
